\origpage{333}{N. Robertson, P.D. Seymour / Journal of Combinatorial Theory, Series B 90 (2004) 325--385}
The purpose of $\mathbf{S}_1$--$\mathbf{S}_4$ is that to prove 3.1, it suffices to prove the following.

\textbf{4.1.} \textit{If $\chi$ satisfies $\mathbf{S}_1$--$\mathbf{S}_4$ then it is not orientedly bad.}

We prove 4.1 in Section 15. The remainder of this section is devoted to proving that 4.1 implies 3.1. To show that, we need several lemmas, which follow.

\textbf{4.2.} \textit{If $\chi$ is a bad colour scheme, there is an orientedly bad colour scheme $\chi'$ with $\Sigma_{\chi'} = \Sigma_\chi$.}

\textbf{Proof.} If $\Sigma_\chi$ is not orientable then $\chi$ itself is orientedly bad. Thus we may assume that $\omega$ is an orientation of $\Sigma_\chi$. Let $(\Gamma_i, \phi_i)$ $(i = 1, 2, \dots)$ be a bad sequence for $\chi$, and let $\Gamma_i = (U_i, V_i, \gamma_i)$ $(i \ge 1)$. Now for each $e \in E(\Gamma_i)$ with $|\tilde{e}| = 3$, let $\gamma'_i(e)$ be a march with $\overline{\gamma'_i(e)} = \tilde{e}$, such that the orientations of $e$ given by $\omega$ and given by the cyclic order of the terms of $\gamma'_i(e)$ coincide, and $\gamma_i(e), \gamma'_i(e)$ have the same first term. For each $e \in E(\Gamma_i)$ with $|\tilde{e}| = 2$ let $\gamma'_i(e) = \gamma_i(e)$. Let $\Gamma'_i = (U_i, V_i, \gamma'_i)$. Let $\phi'_i(e)$ be defined by
\begin{equation*}
\phi'_i(e) = \begin{cases} \phi_i(e) & \text{if } |\tilde{e}| = 2, \\ (\phi_i(e), \pm 1) & \text{if } |\tilde{e}| = 3, \end{cases}
\end{equation*}
where we choose $-1$ if $\gamma_i(e) \ne \gamma'_i(e)$. Let $R$ be the well-quasi-order with $E(R) = \{-1, 1\}$, ordered by equality. Let $\chi'$ be the colour scheme with $\Sigma_{\chi'} = \Sigma_\chi$, $\Phi_{\chi'} = \Phi_\chi$, $\Omega_{\chi'}(3) = \Omega_\chi(3) \times R$, $\Omega_{\chi'}(2) = \Omega_\chi(2)$, and $\Omega_{\chi'}(S) = \Omega_\chi(S)$ for each long side $S$. Then for each $i \ge 1$, $(\Gamma'_i, \phi'_i)$ is a $\chi'$-coloured painting, and $\{(\Gamma'_i, \phi'_i) : i \ge 1\}$ is similarly oriented. Suppose that for some $i, j$ with $j > i \ge 1$, there is a linear inflation $\sigma$ of $(\Gamma'_i, \phi'_i)$ in $(\Gamma'_j, \phi'_j)$. We claim that $\sigma$ is also a linear inflation of $(\Gamma_i, \phi_i)$ in $(\Gamma_j, \phi_j)$. \del{This is mostly clear, but let us} \fix{The paintings $\Gamma_i$ and $\Gamma'_i$ have the same vertices and edges, and differ only in the marches of edges of size 3; and the same holds for $\Gamma_j$ and $\Gamma'_j$. The first two conditions in the definition of an inflation, the first two conditions in the definition of ``linear'', and the definition of ``respects $\Phi_\chi$'' do not involve the marches. So $\sigma$ satisfies them as a map from $\Gamma_i$ to $\Gamma_j$, since it satisfies them as a map from $\Gamma'_i$ to $\Gamma'_j$. For the third condition in the definition of ``linear'', the tail of an edge is the first term of its march, and $\gamma_i(e), \gamma'_i(e)$ have the same first term for every $e \in E(\Gamma_i)$, and similarly for $\Gamma_j$. So each edge has the same tail in $\Gamma_i$ as in $\Gamma'_i$, and in $\Gamma_j$ as in $\Gamma'_j$, and the condition holds for $\sigma$ since it holds for $\sigma$ as a linear inflation of $\Gamma'_i$ in $\Gamma'_j$. Moreover $\phi_i(e) \le \phi_j(\sigma(e))$ for each $e \in E(\Gamma_i)$. If $|\tilde{e}| = 2$ this is $\phi'_i(e) \le \phi'_j(\sigma(e))$; and if $|\tilde{e}| = 3$, then $\sigma(e)$ also has size 3, and it follows by comparing the first terms of $\phi'_i(e)$ and $\phi'_j(\sigma(e))$.}\fixnote{fix:4.2-proof-added}{The paper calls these parts of the claim ``mostly clear''. The argument above is the one used in the formalization.} \fix{It remains to} show that the third condition in the definition of ``inflation'' is satisfied, that is, that for each $e \in E(\Gamma_i)$ and $1 \le k \le |\tilde{e}|$, if $v$ is the $k$th term of $\gamma_i(e)$, then $\sigma(v)$ contains the $k$th term of $\gamma_j(\sigma(e))$. If $\gamma_i(e) = \gamma'_i(e)$ and $\gamma_j(\sigma(e)) = \gamma'_j(\sigma(e))$ then the claim holds since $\sigma$ is an inflation of $\Gamma'_i$ in $\Gamma'_j$. We may assume then that $|\tilde{e}| = 3$. Since the second terms of $\phi'_i(e)$ and of $\phi'_j(\sigma(e))$ are equal, and we may assume they are not both equal to $1$, it follows that they are both $-1$, and consequently $\gamma'_i(e)$ is obtained from $\gamma_i(e)$ by exchanging the second and third terms, and the same holds between $\gamma'_j(\sigma(e))$ and $\gamma_j(\sigma(e))$. If $k = 1$ let $k' = 1$, and if $k = 2$ or $3$ let $k'$ be $3$ or $2$ respectively. Consequently the $k$th term of $\gamma_i(e)$ (that is, $v$) is the $k'$th term of $\gamma'_i(e)$, and the $k$th term of $\gamma_j(\sigma(e))$ is the $k'$th term of $\gamma'_j(\sigma(e))$. Since $\sigma$ is an inflation of $(\Gamma'_i, \phi'_i)$ in $(\Gamma'_j, \phi'_j)$ and $v$ is the $k'$th term of $\gamma'_i(e)$, it follows that $\sigma(v)$ contains the $k'$th term of $\gamma'_j(\sigma(e))$, and therefore contains the $k$th term of $\gamma_j(\sigma(e))$, as required. This proves our claim that $\sigma$ is a linear inflation of $(\Gamma_i, \phi_i)$ in $(\Gamma_j, \phi_j)$, which is impossible since $(\Gamma_i, \phi_i)$ $(i = 1, 2, \dots)$ is a bad sequence for $\chi$. Thus there is no such choice of $i, j, \sigma$, and so $\chi'$ is orientedly bad, as required. $\square$

\textbf{4.3.} \textit{There is no countable sequence $\Omega_i$ $(i = 1, 2, \dots)$ of well-quasi-orders such that $\Omega_{i+1} \prec \Omega_i$ for all $i \ge 1$.}
